\chapter{Shoulder Replacement}
\section*{What Is This Test or Procedure?}
Shoulder arthroplasty replaces damaged parts of the shoulder joint with prosthetic components. Anatomic and reverse designs are selected for different patterns of damage.
\section*{Why This Test or Procedure Is Ordered}
It may be performed for severe arthritis, fracture, cuff-tear arthropathy, or failed prior shoulder reconstruction causing pain and loss of function.
\section*{How This Test or Procedure Is Performed}
The damaged humeral head and sometimes the glenoid are replaced. A reverse replacement changes the joint mechanics when the rotator cuff cannot function adequately.
\section*{Risks and Recovery}
Infection, dislocation, fracture, nerve injury, stiffness, loosening, instability, and revision surgery are possible. A protected period followed by structured shoulder rehabilitation is required.
\section*{References}
\begin{enumerate}\item MedlinePlus. Shoulder Replacement. U.S. National Library of Medicine; 2025.\item Current Medical Diagnosis and Treatment. McGraw-Hill; 2025.\end{enumerate}
\clearpage
